\section{Discussion of the Implementation}

\subsection{Sequential Consistency}
The system provides Sequential Consistency. This means that all clients will
always see the operations in the exact same order, and no one will ever read
"backwards" in time. 

The system achieves this by combining a \textbf{Write Locking
mechanism} (inspired by two-phase locking protocols) with \textbf{Quorum
Intersection ($W + R > N$)} and \textbf{Data Versioning}:

\begin{itemize}
    \item \textbf{Preventing Write Conflicts:} If two clients try to update the
        same key at the exact same time, their Coordinators both send messages
        to the replicas asking for an exclusive lock. Because a write requires
        locking a Quorum (majority) of the nodes, only one Coordinator can win.
        The slower Coordinator receives a "lock denied" message and fails. The
        winner writes the new data with a higher version number. 
    \item \textbf{Preventing Stale Reads:} Because our Write Quorum ($W=2$)
        plus our Read Quorum ($R=2$) is strictly greater than the total number
        of replicas ($N=3$), any read operation must mathematically overlap
        with at least one node that holds the newest successful write. The
        Coordinator simply looks at the version numbers and returns the highest
        one to the client. This makes it impossible for the system to return
        old data after a write has finished.
\end{itemize}

\subsection{Message Efficiency and Data Routing}
The system is designed to use the minimum amount of network messages. We avoid
broadcast messages and keep the back-and-forth communication between nodes as
short as possible.

\begin{itemize}
    \item \textbf{Read and Write Efficiency:} When a client wants to read or
        write, the Coordinator calculates the $N$ responsible nodes and sends
        messages only to them, avoiding network broadcasts. We also minimize
        the round trips. For a read, the Coordinator asks the $N$ nodes and
        they reply with the data—just one round trip, job done. For a write, we
        use exactly two steps (one to lock and read the version, one to save
        the data) to ensure consistency without adding unnecessary confirmation
        messages.
    \item \textbf{Join and Leave Efficiency:} For network changes, we ensure
        data is only sent to the nodes responsible for it. When a node joins,
        it does not ask the whole network for data; it contacts only its
        clockwise neighbor and takes only the keys it needs. When a node
        leaves, it calculates the new network layout to see who will take over
        its keys. It then splits its data and sends it directly to those
        specific nodes, ensuring no data is sent to the wrong places.
\end{itemize}

\subsection{The Write and Lock Workflow}
Because multiple clients can try to update the same data at the same time, we
implemented a strict, two-step locking mechanism. This is the core algorithm
that guarantees our system's consistency. When a Coordinator receives a write
request, it executes the following steps:

\begin{itemize}
    \item \textbf{Step 1: Acquire Locks and Read Version.} 
    The Coordinator sends a \texttt{ReplicaReadRequestMsg} to the $N$
    responsible nodes with a special flag: \texttt{intentToWrite = true}. When
    a replica receives this, it does two things: it locks the specific key so
    no other reads or writes can touch it, and it returns its current data
    version to the Coordinator. 
    
    \item \textbf{Handling Conflicts:} 
    If a second Coordinator tries to write to the same key, the replica sees
    that the key is already locked. It immediately replies with a
    \texttt{lockDenied} message. The second Coordinator fails to reach its
    Quorum and aborts the write. This prevents "dirty writes" and forces
    overlapping operations to happen one after the other.

    \item \textbf{Step 2: Commit and Unlock.} 
    Once the first Coordinator receives successful locks from a Quorum of nodes
    ($W=2$), it finds the highest version number among them, adds 1 to it, and
    sends the \texttt{ReplicaWriteRequestMsg} with the new data. The replicas
    overwrite their old data, save the new version, and finally release the
    lock.

    \item \textbf{Preventing Deadlocks:} 
    A major risk with locks is that a node might stay locked forever if the
    Coordinator crashes in the middle of Step 1. To prevent this, replicas
    start their own internal timeout when they grant a lock. If they do not
    receive the Step 2 commit message within the timeout period, they
    automatically unlock the key, ensuring the system never freezes.
\end{itemize}

\subsection{Additional Assumptions}
To build this system, we made a few specific assumptions about how the hardware
and the network behave:

\begin{itemize}
    \item \textbf{Fail-Stop Crashes:} We assume that when a node fails, it
        simply crashes and stops working. We assume nodes will never act
        maliciously or send corrupted data to the other nodes.
    \item \textbf{Persistent Storage:} We assume that a software crash or power
        loss does not destroy the physical hard drive. When a node recovers, we
        assume its previously saved data is still available on the disk.
    \item \textbf{Unique Node IDs:} We assume the Manager assigns a strictly
        unique ID to every node. No two nodes will ever have the exact same ID
        or occupy the exact same spot on the ring.
\end{itemize}
